% ============================================================
% CASTER-ZT — Supplemental Material: Proofs of Theoretical Properties
% To be submitted as a separate file alongside the main paper.
% ============================================================
\documentclass[lettersize,journal]{IEEEtran}

\usepackage{amsmath,amssymb,amsthm}
\usepackage{cite}
\usepackage{hyperref}

\newtheorem{proposition}{Proposition}

\markboth{IEEE Transactions on Mobile Computing --- Supplemental Material}%
{DK \MakeLowercase{\textit{et al.}}: Proofs of Theoretical Properties}

\begin{document}

\title{Supplemental Material: Proofs of Theoretical Properties\\
for ``Zero-Trust Shielded Graph-Structured Decision Control\\
for Secure Autonomous Recovery in AI-Native 6G\\
Disaster-Monitoring Sensing Networks''}

\author{Georges~Parfait~DK,~Raiwi~El-Rouin,~and~Sonnet~Peter%
\thanks{G.~P.~DK, R.~El-Rouin, and S.~Peter are with the D\'{e}partement de G\'{e}nie Informatique et G\'{e}nie Logiciel, \'{E}cole Polytechnique de Montr\'{e}al, Montr\'{e}al, QC, Canada.}%
}

\maketitle

This supplemental document provides the complete proofs for the ten propositions stated in the main paper.
All proposition numbers and references correspond to those in the main text.

\bigskip

\begin{proof}[Proof of Proposition~1 (Monotone conservatism)]
Since $w_1, w_2, w_3, w_4 > 0$, the composite gate score $g_t = w_1(1-\tau_t)+w_2\rho_t+w_3 u_t+w_4\Delta_t$ is monotonically decreasing in $\tau_t$ and increasing in $\rho_t, u_t, \Delta_t$.
Since $\kappa_1, \kappa_2 \ge 0$, $\tau_{\min} = \tau_0 + \kappa_1\rho_t + \kappa_2 u_t$ is non-decreasing in $\rho_t, u_t$, making the trust condition harder to satisfy.
The \textsc{allow} branch requires $g_t < \gamma_1$, $\tau_t \ge \tau_{\min}$, and $\Delta_t \le \delta_{\max}$.
Worsening any input can only increase $g_t$, raise $\tau_{\min}$, or violate the divergence bound, pushing the output toward a more conservative branch.
\end{proof}

\begin{proof}[Proof of Proposition~2 (Unauthorized-actuation exclusion)]
In the shield algorithm, $\alpha_t=0$ triggers $d_t=\textsc{block}$ at the authorization check, before any permissive branch is evaluated.
\end{proof}

\begin{proof}[Proof of Proposition~3 (Divergence-triggered safeguard)]
$\Delta_t > \delta_{\max}^{\mathrm{hard}}$ triggers $d_t=\textsc{block}$ at the hard-divergence check.
$\Delta_t > \delta_{\max}$ violates the \textsc{allow} condition, preventing permissive decisions.
\end{proof}

\begin{proof}[Proof of Proposition~4 (Bounded decision completeness)]
The shield algorithm is a finite ordered branching structure with mutually exclusive branches (strict threshold ordering) and a final \textsc{else} clause catching all remaining cases.
Every valid input therefore maps to exactly one $d_t \in \mathcal{D}$.
\end{proof}

\begin{proof}[Proof of Proposition~5 (Per-decision complexity)]
Graph encoding: $\mathcal{O}(L|E_t|d)$.
Candidate scoring and trust--risk evaluation: $\mathcal{O}(|\mathcal{A}|k)$ each.
Shield evaluation: $\mathcal{O}(1)$.
Total: $\mathcal{O}(L|E_t|d + |\mathcal{A}|k)$.
\end{proof}

\begin{proof}[Proof of Proposition~6 (Defense-in-depth bound)]
Under identity attack, an unsafe \textsc{allow} requires an authorization false negative: $\Pr \le \varepsilon_\alpha$.
Under telemetry attack: $\Pr \le \varepsilon_\tau$.
Under combined attack with conditional independence:
$\Pr[\textsc{allow} \mid \text{combined}] = \Pr[\alpha_t=1 \mid \text{id}] \cdot \Pr[g_t<\gamma_1 \mid \text{telem}] \le \varepsilon_\alpha \varepsilon_\tau$.
\end{proof}

\begin{proof}[Proof of Proposition~7 (Calibration-consistent accuracy)]
Each episode produces Bernoulli $y_i$ with $\mathbb{E}[y_i]=\varepsilon^\star$.
By Hoeffding's inequality: $\Pr[\varepsilon^\star - \hat\varepsilon > t] \le \exp(-2nt^2)$.
Setting $\exp(-2nt^2)=\delta$ gives $t=\sqrt{\ln(1/\delta)/(2n)}$.
\end{proof}

\begin{proof}[Proof of Proposition~8 (Bounded price of safety)]
With probability $(1-\eta-\eta_{\mathrm{block}})$, the shield allows (loss~0).
With probability $\eta$, scope-reduction yields loss $\le U_{\mathrm{gap}}$.
With probability $\eta_{\mathrm{block}}$, blocking yields loss $\le U_{\max}$.
Summing: $\mathbb{E}[\text{loss}] \le \eta U_{\mathrm{gap}} + \eta_{\mathrm{block}} U_{\max}$.
\end{proof}

\begin{proof}[Proof of Proposition~9 (Conformal admissibility coverage)]
By the conformal framework~\cite{angelopoulos2023conformal}, if scores $(r_1,\ldots,r_{n+1})$ are exchangeable, $\Pr[r_{n+1} \le r_{(\lceil(1-\alpha)(n+1)\rceil)}] \ge 1-\alpha$.
The shield allows only if $g_t < \gamma_1$, which corresponds to $r_{n+1} > r_{(\lceil(1-\alpha)(n+1)\rceil)}$, occurring with probability $\le \alpha$.
\end{proof}

\begin{proof}[Proof of Proposition~10 (Calibration convergence rate)]
By Hoeffding's concentration, $|\hat\varepsilon(\gamma) - \varepsilon(\gamma)| \le \sqrt{\ln(2/\delta)/(2n)}$ w.h.p.
Since $\varepsilon'(\gamma) \ge c_\varepsilon > 0$, the mean value theorem gives $|\varepsilon(\gamma_1^{(n)}) - \varepsilon(\gamma_1^\star)| \ge c_\varepsilon |\gamma_1^{(n)} - \gamma_1^\star|$.
The triangle inequality yields $c_\varepsilon |\gamma_1^{(n)} - \gamma_1^\star| = \mathcal{O}(1/\sqrt{n})$, so $|\gamma_1^{(n)} - \gamma_1^\star| = \mathcal{O}_P(1/\sqrt{n})$.
\end{proof}

\bibliographystyle{IEEEtran}
\bibliography{references}

\end{document}
